
\documentclass{article}
\usepackage{colortbl}
\usepackage{makecell}
\usepackage{multirow}
\usepackage{supertabular}

\begin{document}

\newcounter{utterance}

\centering \large Interaction Transcript for game `eqbench', experiment `eq\_bench', episode 13 with reflexion{-}Qwen3.5{-}4B.
\vspace{24pt}

{ \footnotesize  \setcounter{utterance}{1}
\setlength{\tabcolsep}{0pt}
\begin{supertabular}{c@{$\;$}|p{.15\linewidth}@{}p{.15\linewidth}p{.15\linewidth}p{.15\linewidth}p{.15\linewidth}p{.15\linewidth}}
    \# & \multicolumn{2}{c}{Player} && \multicolumn{2}{c}{Game Master} \\
    \hline

    \theutterance \stepcounter{utterance}  
    & & & \multicolumn{4}{p{0.6\linewidth}}{
        \cellcolor[rgb]{0.9,0.9,0.9}{
            \makecell[{{p{\linewidth}}}]{
                \texttt{\tiny{[P1$\langle$GM]}}
                \texttt{You have attempted this same task before and did not succeed. These are the plans you wrote after those attempts.} \\
\texttt{Trial \#1: New plan: I misinterpreted the emotional arc of the dialogue. Harold's final line ("The same government that can't even balance a budget?") is a direct, personal attack on Lorelei's credibility and logic, not just a statement of fact. This specific rhetorical move is designed to frustrate and annoy the listener, rather than make them feel hopeless (dismay) or resigned (acceptance). While she might feel slightly amused by his cynicism, the primary reaction to being called out for being unrealistic and out of touch is exasperation. I should have recognized that the escalation of Harold's sarcasm and direct challenge to her authority points strongly to a feeling of being worn down and irritated, which aligns with a high score for Exasperation and a near{-}zero score for Dismay. I will ensure my analysis focuses on the specific nature of the final insult to justify the emotional response.} \\
\\ 
\\ 
\\ 
\texttt{Your task is to predict the likely emotional responses of a character in this dialogue:} \\
\\ 
\texttt{Harold: I don't understand why you're so worked up about this election. It's not like the world is going to end if your candidate doesn't win.} \\
\texttt{Lorelei: That's not the point, Harold. It's about what this election represents. It's about the future of our planet, our society.} \\
\texttt{Harold: And you think your candidate has all the answers? } \\
\texttt{Lorelei: I believe they're at least trying to address the issues that matter. Climate change, social justice, healthcare...} \\
\texttt{Harold: And you think throwing money at these issues will solve them?} \\
\texttt{Lorelei: It's not about throwing money, Harold. It's about investing in solutions, in our future.} \\
\texttt{Harold: And who's going to pay for all these grand plans? Us, the taxpayers. } \\
\texttt{Lorelei: Isn't that the point of taxes? To fund public services and infrastructure?} \\
\texttt{Harold: Yes, but there's a limit. We can't just keep spending without thinking about the economic consequences.} \\
\texttt{Lorelei: And we can't keep ignoring the environmental and social consequences of our actions either.} \\
\texttt{Harold: I'm not suggesting we ignore them, Lorelei. I'm just saying we need to be realistic about what we can afford.} \\
\texttt{Lorelei: And I'm saying we can't afford not to act. The cost of inaction is far greater.} \\
\texttt{Harold: And who decides what action is necessary? The government? The same government that can't even balance a budget?} \\
\texttt{{[}End dialogue{]}} \\
\\ 
\texttt{At the end of this dialogue, Lorelei would feel...} \\
\texttt{Dismay} \\
\texttt{Exasperation} \\
\texttt{Amusement} \\
\texttt{Acceptance} \\
\\ 
\texttt{Give each of these possible emotions a score from 0{-}10 for the relative intensity that they are likely to be feeling each.} \\
\\ 
\texttt{You must output in the following format, including headings (of course, you should give your own scores), with no additional commentary:} \\
\\ 
\texttt{Dismay: <score>} \\
\texttt{Exasperation: <score>} \\
\texttt{Amusement: <score>} \\
\texttt{Acceptance: <score>} \\
\\ 
\\ 
\texttt{{[}End of answer{]}} \\
\\ 
\texttt{Remember: zero is a valid score, meaning they are likely not feeling that emotion. You must score at least one emotion > 0.} \\
\\ 
\texttt{Your answer:} \\
            }
        }
    }
     \\ \\

    \theutterance \stepcounter{utterance}  
    & \multicolumn{4}{p{0.6\linewidth}}{
        \cellcolor[rgb]{0.561,0.737,0.561}{
            \makecell[{{p{\linewidth}}}]{
                \texttt{\tiny{[P1$\rangle$GM]}}
                \texttt{Dismay: 0} \\
\texttt{Exasperation: 9} \\
\texttt{Amusement: 2} \\
\texttt{Acceptance: 1} \\
            }
        }
    }
    & & \\ \\

    \theutterance \stepcounter{utterance}  
    & & & \multicolumn{2}{p{0.3\linewidth}}{
        \cellcolor[rgb]{0.9,0.9,0.9}{
            \makecell[{{p{\linewidth}}}]{
                \texttt{\tiny{[GM$|$GM]}}
                \texttt{Dismay: 0} \\
\texttt{Exasperation: 9} \\
\texttt{Amusement: 2} \\
\texttt{Acceptance: 1} \\
            }
        }
    }
    & & \\ \\

    \theutterance \stepcounter{utterance}  
    & & & \multicolumn{2}{p{0.3\linewidth}}{
        \cellcolor[rgb]{0.9,0.9,0.9}{
            \makecell[{{p{\linewidth}}}]{
                \texttt{\tiny{[GM$|$GM]}}
                \texttt{Dismay: 4} \\
\texttt{Exasperation: 8} \\
\texttt{Amusement: 0} \\
\texttt{Acceptance: 0} \\
            }
        }
    }
    & & \\ \\

    \theutterance \stepcounter{utterance}  
    & & & \multicolumn{2}{p{0.3\linewidth}}{
        \cellcolor[rgb]{0.9,0.9,0.9}{
            \makecell[{{p{\linewidth}}}]{
                \texttt{\tiny{[GM$|$GM]}}
                \texttt{game\_result = LOSE} \\
            }
        }
    }
    & & \\ \\

    \theutterance \stepcounter{utterance}  
    & & \multicolumn{4}{p{0.6\linewidth}}{
        \cellcolor[rgb]{0.95,0.95,0.95}{
            \makecell[{{p{\linewidth}}}]{
                \texttt{\tiny{[P1$\leftrightarrow$P]}}
                \texttt{Post{-}episode Reflexion (not sent to GM):} \\
\texttt{New plan: I correctly identified Exasperation as the primary emotion, but I likely underestimated the intensity of Dismay. Harold's final line ("The same government that can't even balance a budget?") is a devastating rhetorical blow that undermines the very foundation of Lorelei's argument (that the government can act). This doesn't just annoy her; it makes her feel hopeless and defeated by the reality of the situation, which is the definition of dismay. While she is exasperated by the attack, the content of the attack (the impossibility of the task) induces a deeper sense of despair. I will adjust the score for Dismay to be high (e.g., 8 or 9) to reflect this crushing realization, while keeping Exasperation high as the immediate reaction to the insult. I will ensure the scores reflect a progression from irritation (exasperation) to hopelessness (dismay) caused by the specific content of the final insult.} \\
            }
        }
    }
    & \\ \\

\end{supertabular}
}

\end{document}
